\section{Further questions}\label{sec:questions}

The preceding theorems distinguish three kinds of information.
The transported energy controls the complete weak linearized
spectrum. The nonlinear reconstruction also requires the
component and collar estimates for the actual Chern curvature.
Integral monodromy is obtained only after the marked observations
have determined the missing incidence and the first-order
smoothing has fixed its integral pairing. Several extensions
require additional arguments at precisely these distinctions.

\subsection{The weakest nonlinear time}
Proposition~\ref{prop:three-scales} gives a compact stable
rank-two family whose weakest eigenvalue has order $h^4$,
whereas its largest edge weight has order $h$.
Thus $L_*\asymp h^{-1}$ and $L_*\nu_0\asymp h^3$.
Theorem~\ref{thm:tracking} controls bounded intervals
of order $h^{-1}$. Its residual estimate, integrated
to order $h^{-4}$, gives a bound of order
$h^{-2}(1+|\log h|)^2$, which does not tend to zero.
The all-time hypothesis of Theorem~\ref{thm:all-time}
therefore cannot be removed by the same argument.

\begin{question}
For the family of Proposition~\ref{prop:three-scales},
does a finite nonlinear evolution track the actual
HYM flow uniformly through the time $h^{-4}$ for
the spatial initial class of Section~\ref{sec:entry}?
If so, is the leading equation
\eqref{eq:finite-flow} sufficient, or must it be
corrected?
\end{question}

Failure of the integrated upper bound is not a
counterexample to tracking. Both exact and finite
solutions can retain a weak component for a long
time. A useful proof might estimate the residual
relative to the evolving dissipation, or incorporate
the next terms in the finite field. Either approach
must retain cancellations among the incident
forces; a bound on each force separately can lose
the scale of their sum. Relative linear spectral
comparison alone does not settle this question.

\subsection{Independent nodes with component holonomy}
Section~\ref{sec:component-holonomy} permits arbitrary
ratios of the component energy to the common neck
energy. Sections~\ref{sec:geometry}--\ref{sec:entry}
permit arbitrary ratios of the six node energies
when the component connections are trivial.
Combining these results requires a compact family
with both sets of independent parameters and an
analysis on that family. In particular, the
harmonic gauges on the punctured components must
give energy-controlled boundary estimates at each
individual neck, and the nonlinear matching must
retain the component connection terms. An abstract
sum of the two proved operators would not verify
these requirements.

\subsection{Nonunitary degeneration and singular limits}
All bundles used in this paper admit flat unitary
reference connections. Their holomorphic structures
are polystable, and the degenerations of stability
studied here concern reducibility and the
appearance of parallel endomorphisms. Genuine
nonunitary extensions can introduce additional
deformation and obstruction terms into the
effective equation. Unstable holomorphic bundles
also lack the stationary reference metric used
in the order bounds and tail estimates.
Their actual HYM evolution requires a different
normalization and confinement argument.

On a smooth curve, varying the Kähler form alone
does not produce ordinary slope walls. A broader
stability problem must be realized by varying
holomorphic bundle data, a specified parabolic
stability structure, or higher-dimensional
polarizations. When singular graded sheaves or
curvature concentration occur, their local charge
and scale data must enter the matching problem.
No singular concentration law follows from the
flat-bundle results proved here.

\subsection{Conditioning and additional observations}
For a known graph with independent weights,
Theorem~\ref{thm:projective-inverse} resolves an edge
only when the observed block error is small
relative to that edge's weight. An observation
scheme adapted to successively weaker spectral
subspaces might recover smaller weights, but
would require compatible markings and a sharper
operator estimate. Eigenvalue accuracy by itself
does not provide these identifications.

For unknown incidence,
Theorem~\ref{thm:monodromy-inverse} uses a finite
simple-graph class, known first-order node
smoothing and component averages of eigenvectors.
Parallel edges, unknown smoothing multiplicities
and unmarked component permutations require
different observations. It would be useful to
identify a minimal sufficient observation package
and to quantify the propagation of errors from
raw eigenvalue and eigensection measurements.
The present operator-level error margin does not
answer that statistical or numerical question.

\subsection{Higher-dimensional geometry}
The curve construction separates component
parallel endomorphisms from the directions
coupled by the nodes. In higher dimensions
both that decomposition and the local capacity
law can change. A bundle theorem must identify
the actual endomorphism complex, construct its
compact holomorphic family, and prove the
uniform analytic estimates in the chosen
Kähler metrics. Small eigenvalues of ordinary
differential forms do not automatically give
the required endomorphism modes.
Likewise, recovering projective bundle transport
does not recover the integral cohomology of
the base. A higher-dimensional synthesis must
specify and prove each connection on the same
geometric family.
